\documentclass[11pt,a4paper]{article}

\usepackage[utf8]{inputenc}
\usepackage[T1]{fontenc}
\usepackage[spanish]{babel}
\usepackage{amsmath,amssymb,mathtools}
\usepackage{geometry}
\usepackage{lmodern}
\usepackage{microtype}

\geometry{margin=2.4cm}
\setlength{\parindent}{0pt}
\setlength{\parskip}{0.45em}

\newcommand{\R}{\mathbb{R}}
\newcommand{\sen}{\operatorname{sen}}

\begin{document}

\begin{center}
{\LARGE\bfseries Análisis Matemático III}\\[0.4em]
{\Large Tema 1: Conceptos básicos de cálculo vectorial}\\[0.4em]
{\large Enunciados}
\end{center}

\section{Ejercicio 1}


\section{Ejercicio 2}



\section{Ejercicio 3}



\section{Ejercicio 4}

Determina los puntos dobles de la curva

\[
x^3+x^2+y^2-x-4y+3=0
\]

y analiza de qué tipo son.

\section{Ejercicio 5}

Determina los puntos dobles de la curva

\[
x^3+y^3=3axy,
\qquad
a\neq0
\]

y analiza de qué tipo son.

\section{Ejercicio 6}

Determina los puntos dobles de la curva

\[
y^3=x^3+ax^2,
\qquad
a\neq0
\]

y analiza de qué tipo son.

\section{Ejercicio 7}

Determina los puntos dobles de la curva

\[
y^2=2x^2y+x^4y-2x^4
\]

y analiza de qué tipo son.

\section{Ejercicio 8}

Determina si las siguientes parametrizaciones son regulares o no. En caso de que no lo sean, comprueba si pudiera existir una parametrización regular de la misma curva. Finalmente, representa gráficamente las curvas asociadas:

\[
\text{a) }\alpha(t)=(t^3,t^6)
\]

\[
\text{b) }\beta(t)=(t^2,t^3)
\]

\[
\text{c) }\gamma(t)=(\cos^3(t),\sen^3(t)),
\qquad
t\in[0,2\pi].
\]

\section{Ejercicio 9}

Determina los puntos de la curva dada por la parametrización

\[
\alpha(t)=(t^3-3t,t^2-4)
\]

en los que su tangente sea vertical u horizontal.

\section{Ejercicio 10}

Obtén la expresión de la recta tangente a la circunferencia unidad centrada en el origen en cualquier punto.

\section{Ejercicio 11}

Determina las rectas tangente y normal a la curva

\[
y=x^3
\]

en el punto

\[
(1,1)
\]

utilizando su expresión paramétrica.

\section{Ejercicio 12}

Calcula mediante la fórmula integral paramétrica la longitud de la circunferencia de radio

\[
r=3
\]

que está centrada en el punto

\[
(-1,1).
\]

\section{Ejercicio 13}

Calcula la longitud de la curva dada por la función

\[
y=\sqrt{x^3}
\]

para

\[
x\in[0,3].
\]

\section{Ejercicio 14}

Determina la representación paramétrica de la curva dada por la intersección de la esfera centrada en el origen y radio \(2\) y el plano

\[
y+z=2.
\]

\section{Ejercicio 15}

Determina la representación paramétrica de la curva dada por la intersección de la semiesfera superior centrada en el origen y radio \(2\) y el cilindro dado por la ecuación

\[
(x-1)^2+y^2=1,
\]

también conocida como la curva de Viviani.

\section{Ejercicio 16}

Dada la curva \(C\) con parametrización

\[
\alpha(t)
=
\left(
t^2-9,\frac{t^3}{3},2t
\right)
\]

y los puntos

\[
P=(-9,0,0),
\qquad
Q=(0,9,6),
\]

calcula la longitud del arco de curva entre los puntos \(P\) y \(Q\).

\section{Ejercicio 17}

Se considera la curva \(C\) con parametrización

\[
\alpha(t)
=
(t-\sen(t),t,1-\cos(t)),
\qquad
t\in[0,6\pi].
\]

\[
\text{a) Determina si la parametrización es regular.}
\]

\[
\text{b) Plantea la integral que habría que calcular para obtener la longitud de la curva.}
\]

\section{Ejercicio 18}

Determina la expresión de la recta tangente a la gráfica de la curva dada por la parametrización

\[
\alpha(t)
=
(t^2+1,2+\ln(t+1),e^t)
\]

en el punto

\[
(1,2,1).
\]

\section{Ejercicio 19}

Calcula los vectores velocidad y aceleración, así como la celeridad, de la curva dada por la parametrización

\[
\alpha(t)
=
(t^3,1-t,4t^2)
\]

en el punto asociado al valor

\[
t=1.
\]

\section{Ejercicio 20}

Dada la curva con parametrización

\[
\alpha(t)
=
(t^2,1+t,1+t^3),
\]

comprueba que la curva pasa por el punto

\[
(1,0,0)
\]

y calcula los vectores velocidad y aceleración en dicho punto.

\section{Ejercicio 21}

Dada una partícula de la que se conoce su vector aceleración

\[
a(t)
=
(e^t,0,t+1)
\]

y su velocidad inicial

\[
v(0)
=
(1,-3,\sqrt2),
\]

calcula su vector velocidad.

\section{Ejercicio 22}

Determina si la curva de parametrización

\[
\alpha(t)
=
(t,2t,-3t)
\]

incide perpendicularmente en el plano

\[
x+2y-3z=10,
\]

y si es así calcula el punto de intersección.

\section{Ejercicio 23}

De una curva con parametrización \(\alpha(t)\) se sabe que

\[
\alpha''(t)
=
(0,0,2)
\qquad
\forall\,t\in\R,
\]

\[
\alpha'(0)
=
(2,0,0)
\]

y

\[
\alpha(0)
=
(2,0,0).
\]

A partir de esos datos calcula \(\alpha(t)\) y el vector tangente unitario en

\[
t=1.
\]

\section{Ejercicio 24}

Dada la curva con parametrización

\[
\alpha(t)
=
(2\cos(t),2\sen(t),5t),
\]

calcula la recta tangente y el plano normal a la curva en el punto

\[
\left(
0,2,\frac{5\pi}{2}
\right).
\]

\section{Ejercicio 25}

Descompón el vector aceleración de la curva con parametrización

\[
\alpha(t)
=
(t,t^2,t)
\]

en el punto asociado al valor

\[
t=1
\]

en sus componentes normal y tangencial.

\section{Ejercicio 26}

Determina si los conjuntos

\[
D_1
=
\R^2\setminus\{(0,0)\}
\]

y

\[
D_2
=
\R^3\setminus\{(0,0,0)\}
\]

son simplemente conexos.

\section{Ejercicio 27}

Determina si las superficies asociadas a una esfera y a un toro son simplemente conexas.

\section{Ejercicio 28}

Determina si

\[
D_1
=
\left\{
(x,y,z)\in\R^3
\mid
1<x^2+y^2+z^2<4
\right\}
\]

y

\[
D_2
=
\left\{
(x,y,z)\in\R^3
\mid
x^2+y^2\neq0
\right\}
\]

son conjuntos simplemente conexos.

\section{Ejercicio 29}

Se considera la superficie \(S\) con representación paramétrica

\[
\vec r(u,v)
=
(u,v,u^2+\sen(v)),
\]

con

\[
u\in[-2,2]
\]

y

\[
v\in[0,2\pi].
\]

Determina si la parametrización es regular.

\section{Ejercicio 30}

Dada la superficie \(S\) de representación paramétrica

\[
\vec r(u,v)
=
(u+v^2,u+4v,u^2-u),
\]

con

\[
(u,v)\in\R^2,
\]

determina si la parametrización es regular.

\end{document}
